\chapter{Results}
\label{chap:resu}

This chapter presents the experimental results of the intrusion detection system deployed on real hardware, organized around the three research questions formulated in Chapter~\ref{chap:Introduction}. First, the performance of classical ML models is analyzed as a centralized baseline (Section~\ref{sec:res_classic}). Then, the hierarchical federated learning system is evaluated (Section~\ref{sec:res_hfl}). Finally, the overhead of ASCON-128 encryption is quantified (Section~\ref{sec:res_ascon}).

% ====================================================================
\section{Classical Model Performance (Centralized Baseline)}
\label{sec:res_classic}

The four classical ML models were trained on the entire dataset (256,276 flows) with stratified 80/20 partitioning and hyperparameter optimization via cross-validation ($k=5$). Table~\ref{tab:classic_results} summarizes the results obtained.

\begin{table}[H]
\centering
\caption{Classical ML model performance (centralized training).}
\label{tab:classic_results}
\small
\begin{tabular}{l c c c c c c}
\toprule
\textbf{Model} & \textbf{Accuracy} & \textbf{F1 Macro} & \textbf{F1 Weighted} & \textbf{Precision} & \textbf{Recall} & \textbf{MCC} \\
\midrule
Decision Tree      & 0.9987 & 0.9981 & 0.9987 & 0.9982 & 0.9981 & 0.9977 \\
Random Forest      & 0.9994 & 0.9991 & 0.9994 & 0.9991 & 0.9990 & 0.9989 \\
Logistic Regr.     & 0.9412 & 0.9118 & 0.9385 & 0.9244 & 0.9053 & 0.8964 \\
SVM (Linear)       & 0.9447 & 0.9162 & 0.9421 & 0.9280 & 0.9099 & 0.9017 \\
SVM (RBF)          & 0.9968 & 0.9953 & 0.9968 & 0.9954 & 0.9952 & 0.9944 \\
\bottomrule
\end{tabular}
\end{table}

\textbf{Key findings:}
\begin{itemize}
    \item \textbf{Random Forest} achieves the best overall performance (Accuracy: 99.94\%, MCC: 0.9989), followed closely by Decision Tree (99.87\%) and SVM-RBF (99.68\%).
    \item Linear models (Logistic Regression, SVM-Linear) show a significant drop ($\approx 5$--$6$ percentage points), suggesting that decision boundaries between classes are not completely linearly separable in the 13-feature space.
    \item All tree-based models exhibit MCC $> 0.99$, confirming that the 13 selected features are highly discriminative for the three-class classification task.
\end{itemize}

\subsection{ESP32 Deployment Feasibility}

Table~\ref{tab:classic_esp32} evaluates the feasibility of exporting each model to C code for execution on the ESP32-S3.

\begin{table}[H]
\centering
\caption{Estimated memory footprint for ESP32 deployment.}
\label{tab:classic_esp32}
\small
\begin{tabular}{l r r l}
\toprule
\textbf{Model} & \textbf{Est. Memory} & \textbf{Inference (ms)} & \textbf{ESP32 Viable?} \\
\midrule
Decision Tree       & $\sim$15\,KB  & $< 0.1$ & Yes \\
Random Forest       & $\sim$2.5\,MB & $\sim$5  & No (exceeds SRAM) \\
Logistic Regression & $\sim$1\,KB   & $< 0.05$ & Yes \\
SVM (Linear)        & $\sim$1\,KB   & $< 0.05$ & Yes \\
SVM (RBF)           & $\sim$4\,MB   & $\sim$50 & No (exceeds SRAM) \\
MLP (proposed)      & $\sim$4.5\,KB & $< 0.1$  & \textbf{Yes} \\
\bottomrule
\end{tabular}
\end{table}

This reveals a fundamental \emph{trade-off}: the most accurate models (Random Forest, SVM-RBF) are incompatible with the ESP32-S3's memory constraints (512\,KB SRAM). The proposed MLP ($\approx 4.5$\,KB) offers an efficient compromise between accuracy and edge viability.

% ====================================================================
\section{Hierarchical Federated Learning: Convergence and Accuracy}
\label{sec:res_hfl}

The HFL system was deployed on real hardware during sessions of $\geq 60$ federated learning rounds.

\subsection{Federated Model Convergence}

\begin{table}[H]
\centering
\caption{HFL system convergence metrics.}
\label{tab:hfl_convergence}
\begin{tabular}{l r}
\toprule
\textbf{Metric} & \textbf{Value} \\
\midrule
Rounds to accuracy $> 70\%$  & $\sim$8--12 \\
Rounds to accuracy $> 85\%$  & $\sim$20--30 \\
Stabilized global accuracy    & $\sim$88--92\% \\
Final loss (Cross-Entropy)    & $\sim$0.25--0.40 \\
Samples per round (per gateway) & 40 \\
Local epochs per round        & 5 \\
Duration per round            & $\sim$70--80\,s \\
\bottomrule
\end{tabular}
\end{table}

\textbf{Analysis:}
\begin{enumerate}
    \item The federated model converges in $\sim$20--30 rounds to $\sim$90\% accuracy, representing a \textbf{gap of $\sim$10 points} relative to centralized models (Decision Tree: 99.87\%). This gap is attributable to: (i) the reduced local buffer size (40 samples vs.\ 256,276), (ii) the non-IID data distribution between nodes (one node generates only normal traffic), and (iii) the heuristic labeling function that introduces noise relative to real labels.

    \item Convergence is monotonically increasing during the first 20 rounds, with minor oscillations ($\pm 3\%$) in subsequent rounds, typical of FL with non-IID data \cite{Imteaj2022_FL_ResourceConstrainedSurvey}.

    \item The \textbf{communication cost} per complete federated round (round trip) is $\approx 7.2$\,KB in plain JSON and $\approx 9.8$\,KB with ASCON encryption, which is compatible with low-bandwidth IoT networks.
\end{enumerate}

\subsection{Comparison: Federated vs.\ Centralized}

Table~\ref{tab:fed_vs_central} directly compares federated model accuracy with centralized models, considering that the federated model uses an identical MLP.

\begin{table}[H]
\centering
\caption{Accuracy comparison: federated vs.\ centralized training.}
\label{tab:fed_vs_central}
\small
\begin{tabular}{l c c l}
\toprule
\textbf{Approach} & \textbf{Accuracy} & \textbf{Data Shared} & \textbf{Privacy} \\
\midrule
Centralized MLP              & $\sim$97\%    & All flows             & None \\
Federated MLP (HFL, 60 rounds) & $\sim$90\%  & Weights only (163 params) & High \\
Centralized Random Forest    & 99.94\%       & All flows             & None \\
Centralized Decision Tree    & 99.87\%       & All flows             & None \\
\bottomrule
\end{tabular}
\end{table}

The federated system achieves $\sim$90\% accuracy sharing \textbf{only 163 model parameters} (652 bytes) instead of raw traffic data ($\sim$256K flows $\times$ 13 features $= \sim$13.3M values). This represents a \textbf{$>$99.99\% reduction} in exchanged information, with a $\sim$7--10 point accuracy penalty.

% ====================================================================
\section{ASCON-128 Encryption Overhead}
\label{sec:res_ascon}

Encryption/decryption metrics were collected during 60 federated learning rounds (182 recorded cryptographic operations). Results are segmented by communication direction and system layer.

\subsection{Temporal Overhead}

\begin{table}[H]
\centering
\caption{ASCON-128 encryption/decryption times on the central server (PC).}
\label{tab:ascon_time}
\small
\begin{tabular}{l c c c c r}
\toprule
\textbf{Operation} & \textbf{Mean} & \textbf{Median} & \textbf{Min} & \textbf{Max} & \textbf{N} \\
\midrule
Decryption (RPi$\to$PC) & 14.21\,ms & 12.69\,ms & 10.62\,ms & 43.73\,ms & 122 \\
Encryption (PC$\to$RPi)  & 13.76\,ms & 12.64\,ms & 10.10\,ms & 93.42\,ms &  60 \\
\midrule
\textbf{Total (average)} & \textbf{13.98\,ms} & \textbf{12.67\,ms} & --- & --- & 182 \\
\bottomrule
\end{tabular}
\end{table}

\textbf{Observations:}
\begin{itemize}
    \item The average time for an ASCON-128 operation in Python (PC server) is \textbf{$\approx$14\,ms}, with a median of $\approx$12.7\,ms. Given that each FL round takes $\sim$70--80 seconds total, the cryptographic overhead per round (3 operations: 2 decryptions + 1 encryption) represents $\sim$42\,ms, equivalent to \textbf{$\approx$0.05\%} of total round time.
    \item Outliers (max 93.42\,ms) correspond to occasional OS load spikes and are not representative of typical behavior.
    \item On the ESP32 (C++, measured with \texttt{micros()}), encryption/decryption times for feature payloads ($\sim$200 bytes) are on the order of \textbf{$\sim$100--300\,$\mu$s}, thanks to native C implementation without interpreter overhead.
\end{itemize}

\subsection{Communication Overhead}

\begin{table}[H]
\centering
\caption{Size overhead from ASCON-128 + Base64 encryption.}
\label{tab:ascon_size}
\small
\begin{tabular}{l r r r r}
\toprule
\textbf{Channel} & \textbf{Plaintext} & \textbf{Encrypted} & \textbf{Overhead} & \textbf{Increase} \\
\midrule
RPi$\to$PC (weights)       & 3,636$\pm$9\,B  & 4,948$\pm$14\,B  & 1,312$\pm$4\,B  & 36.1\% \\
PC$\to$RPi (global model)  & 3,530$\pm$8\,B  & 4,790$\pm$11\,B  & 1,260$\pm$3\,B  & 35.7\% \\
\midrule
\textbf{Average}           & 3,583\,B         & 4,869\,B          & \textbf{1,286\,B} & \textbf{35.9\%} \\
\bottomrule
\end{tabular}
\end{table}

The $\approx$36\% communication overhead is composed of:
\begin{itemize}
    \item \textbf{ASCON authentication tag:} 16 bytes (fixed).
    \item \textbf{Nonce:} 16 bytes (fixed).
    \item \textbf{Base64 expansion:} factor $\times 4/3$ over ciphertext + tag + nonce.
    \item \textbf{JSON envelope:} $\sim$30 bytes of structure (\texttt{\{"ct":"...","tag":"...","nonce":"..."\}}).
\end{itemize}

In absolute terms, the overhead is $\approx$1.3\,KB per message, which for the $\sim$3 transmissions per FL round amounts to $\sim$3.9\,KB additional per round. For a standard Wi-Fi link (54\,Mbps), this adds $<$0.001\,ms of transmission time, making it negligible.

\subsection{Summary: ASCON Impact on the HFL System}

\begin{table}[H]
\centering
\caption{Total ASCON-128 overhead summary per federated round.}
\label{tab:ascon_resumen}
\begin{tabular}{l r r}
\toprule
\textbf{Component} & \textbf{With ASCON} & \textbf{Without ASCON} \\
\midrule
Communication time (PC, 3 ops.)  & $\sim$42\,ms   & $\sim$0.3\,ms (serialization) \\
Total size per round              & $\sim$14.6\,KB & $\sim$10.7\,KB \\
Total round time                  & $\sim$73\,s    & $\sim$72\,s \\
\midrule
\textbf{Relative temporal overhead}  & \multicolumn{2}{c}{\textbf{$\approx$0.06\% of FL cycle}} \\
\textbf{Relative size overhead}      & \multicolumn{2}{c}{\textbf{$\approx$36\% per message}} \\
\bottomrule
\end{tabular}
\end{table}

ASCON-128 encryption introduces a \textbf{negligible} temporal overhead ($< 0.1\%$ of total round time) and a \textbf{moderate} size overhead ($\approx 36\%$) resulting in $\sim$3.9\,KB additional per round. Given that the payload is on the order of kilobytes, this cost is acceptable for typical IoT networks and does not impact federated model convergence or performance.

% ====================================================================
\section{Research Questions Analysis}
\label{sec:preguntas}

\subsection{RQ-1: Can TinyML detect network attacks on IoT devices with accuracy comparable to traditional edge models?}

The 1,139-parameter MLP model ($\sim$4.5\,KB), executed directly on the ESP32-S3, achieves $\sim$90\% accuracy after federated training. While a gap of $\sim$10 points exists relative to centralized Random Forest (99.94\%), the latter requires $\sim$2.5\,MB of memory and cannot be deployed on a microcontroller with 512\,KB SRAM.

Comparing with edge-viable models (Decision Tree: 99.87\% with $\sim$15\,KB), the gap narrows to $\sim$10 points, primarily attributable to the federated scheme rather than the model architecture. An identical MLP trained centrally achieves $\sim$97\%, placing it within 3 points of the Decision Tree.

\textbf{RQ-1 Conclusion:} TinyML can detect attacks with production-viable accuracy ($>$90\%) in a form factor compatible with IoT microcontrollers. The penalty relative to traditional edge models is moderate and compensated by privacy (FL) and security (ASCON) advantages.

\subsection{RQ-2: Does federated learning reduce sensitive data exchange without affecting IDS accuracy?}

Federated learning reduces exchanged information from $\sim$13.3 million values (all flows) to \textbf{163 parameters} per round (652\,bytes), a reduction of $>$99.99\%. The accuracy impact is $\sim$7 points (from $\sim$97\% centralized to $\sim$90\% federated), consistent with FL literature on non-IID data \cite{Imteaj2022_FL_ResourceConstrainedSurvey, Albanbay2025}.

\textbf{RQ-2 Conclusion:} Federated learning eliminates the need to transfer raw traffic data, protecting monitored network privacy. The accuracy reduction ($\sim$7\%) is an acceptable trade-off in scenarios where privacy is a priority.

\subsection{RQ-3: What is the overhead of ASCON encryption in federated model communication?}

Based on 182 cryptographic operations recorded during 60 FL rounds:
\begin{itemize}
    \item \textbf{Temporal overhead:} $\sim$42\,ms per round ($\sim$14\,ms per operation), representing $\sim$0.06\% of total round time.
    \item \textbf{Size overhead:} $\sim$36\% increase per message ($\sim$1.3\,KB over $\sim$3.6\,KB).
    \item \textbf{Convergence impact:} None observable. Accuracy and loss curves are indistinguishable between ASCON-enabled and ASCON-disabled branches.
\end{itemize}

\textbf{RQ-3 Conclusion:} ASCON-128 overhead is \textbf{negligible in time} and \textbf{moderate in size}, validating its suitability as a protection mechanism for federated learning communications on constrained IoT devices.
